% Reference exam skeleton — hand-author from a question bank, then build with reg-exam-build --roster. New format per notes/exam-tex-doctrine.
%
% Usage:
%   1. Copy this file to your exam directory and rename it (e.g. cecs378-midterm1-fa26.tex).
%   2. Edit the \metastrip{}{}{}{} arguments: course, term, exam name, total points.
%   3. Edit the running header (\fancyhead lines) to match course/term/exam.
%   4. Replace the placeholder sections below with your actual questions.
%   5. Build: reg-exam-build --roster roster.csv cecs378-midterm1-fa26.tex
%
% A.2 per-student ID macros are in the preamble — do not remove them.
% The \examserial, \studentname, \studentid, and \studentserial macros are all
% injected at build time by reg-exam-build; the \@ifundefined+\def defaults here
% are the fallback for stand-alone pdflatex compiles.

\documentclass[11pt]{article}
\usepackage[margin=0.9in,top=0.7in,bottom=0.85in,headheight=14pt,headsep=14pt]{geometry}
\usepackage[T1]{fontenc}
\usepackage[utf8]{inputenc}

% --- Print-optimized font stack ---
\usepackage{charter}                          % Bitstream Charter body
\usepackage[scaled=0.92]{helvet}              % Helvetica-clone sans (metadata)
\usepackage[scaled=0.95,varl]{inconsolata}    % Inconsolata mono, varl alt
\renewcommand{\familydefault}{\rmdefault}
\linespread{1.06}
\frenchspacing

\usepackage{enumitem}
\usepackage{needspace}
\usepackage[hyphens]{url}
\usepackage[hidelinks]{hyperref}
\usepackage{listings}
\usepackage{microtype}
\usepackage{fancyhdr}
\usepackage{pifont}                           % Zapf Dingbats — manicule (\ding{43})
\usepackage{xcolor}
\definecolor{keyred}{HTML}{B71C1C}
\definecolor{keygreen}{HTML}{1B5E20}
\definecolor{keygrey}{HTML}{616161}

% MC choice color: green-bold for correct, grey for distractors. Gated by \ifanswers
% so student PDFs render identically to the pre-color source.
\newcommand{\correctchoice}[1]{\ifanswers\textcolor{keygreen}{\textbf{#1}}\else #1\fi}
\newcommand{\wrongchoice}[1]{\ifanswers\textcolor{keygrey}{#1}\else #1\fi}

% --- Build-time stamp (pdfTeX primitives) ---
\newcount\stamphour
\newcount\stampminute
\stamphour=\time
\divide\stamphour by 60
\stampminute=\time
\count255=\stamphour
\multiply\count255 by 60
\advance\stampminute by -\count255
\newcommand{\padtwo}[1]{\ifnum#1<10 0\fi\the#1}
\newcommand{\buildtime}{\the\year-\padtwo{\month}-\padtwo{\day}~\padtwo{\stamphour}:\padtwo{\stampminute}}

% --- Answer key toggle ---
% reg-exam-build injects \def\answersmode{1} for the key build; bridge that
% to \ifanswers so the key PDF reveals answers.
\newif\ifanswers
\answersfalse
\makeatletter
\@ifundefined{answersmode}{}{\answerstrue}
\makeatother

% --- Exam serial (set at build time by reg-exam-build) ---
\providecommand{\examserial}{XXXXXXXX}

% --- Per-student exam ID macros (A.2 doctrine) ---
% reg-exam-build injects \def\studentname, \def\studentid, \def\studentserial
% before \input. These defaults are the stand-alone compile fallback only.
% NB: use \@ifundefined+\def (not \providecommand) so the empty default is a
% non-\long macro that \ifx-compares equal to \@empty — the prefill guards below
% rely on it. \providecommand{}{} makes a \long macro that is NOT \ifx-\@empty.
\makeatletter
\@ifundefined{studentname}{\def\studentname{}}{}
\@ifundefined{studentid}{\def\studentid{}}{}
\@ifundefined{studentserial}{\def\studentserial{}}{}
\makeatother

% \fieldline{width}{value} — value rests ON the rule (empty => blank line for the
% student to write in). \rlap gives the value zero width so the rule draws full
% width; \raisebox lifts the baseline so text sits on the line, not through it.
\newcommand{\fieldline}[2]{%
  \rlap{\hspace{0.5em}\raisebox{2.2pt}{#2}}\rule[-3pt]{#1}{0.8pt}%
}
% \identityinstruction — pre-filled exams ask the student to VERIFY; blank copies
% keep the PRINT-CLEARLY warning.
\makeatletter
\newcommand{\identityinstruction}{%
  \ifx\studentname\@empty
    {\sffamily\bfseries\small\ding{43}~~PRINT CLEARLY. UNNAMED EXAMS CANNOT BE RETURNED OR GRADED.}%
  \else
    {\sffamily\bfseries\small\ding{43}~~VERIFY YOUR NAME AND STUDENT ID BELOW ARE CORRECT, THEN ADD TODAY'S DATE.}%
  \fi%
}
\makeatother

% --- Running header + italic foot ---
\pagestyle{fancy}
\fancyhf{}
\fancyhead[L]{\textit{\small CECS NNN · Term YYYY}}   % <-- edit course/term
\fancyhead[R]{\textit{\small Exam Name}}               % <-- edit exam name
\fancyfoot[L]{\textit{\footnotesize Generated \buildtime}}
\fancyfoot[C]{\textit{\small p.~\thepage}}
\makeatletter
\fancyfoot[R]{\textit{\footnotesize Serial \texttt{\examserial}%
  \ifx\studentserial\@empty\else
    \ \textperiodcentered\ ID \texttt{\studentserial}%
  \fi}}
\makeatother
\renewcommand{\headrulewidth}{0.4pt}
\renewcommand{\footrulewidth}{0pt}

% --- Letterpress double rule ---
\newcommand{\doublerule}{%
  \noindent\rule{\linewidth}{1.0pt}\par\vspace{1.5pt}%
  \noindent\rule{\linewidth}{0.3pt}\par}

% --- Drafting-style metadata strip ---
\newcommand{\metastrip}[4]{%
  \par\noindent\rule{\linewidth}{0.3pt}\par\vspace{2pt}%
  \noindent\begin{minipage}{\linewidth}%
  \begin{tabbing}%
  \hspace{6em}\=\hspace{16em}\=\hspace{6em}\=\kill%
  \textsc{course}~~\>#1\>\textsc{term}~~\>#2\\[1pt]%
  \textsc{exam}~~\>#3\>\textsc{points}~~\>#4%
  \end{tabbing}%
  \end{minipage}\par\vspace{-4pt}%
  \noindent\rule{\linewidth}{0.3pt}\par}

% --- Section heading ---
\newcommand{\examsection}[1]{%
  \par\vspace{1.4em}%
  \noindent\rule{\linewidth}{0.3pt}\par\vspace{2pt}%
  \noindent{\textsc{\large #1}}%
  \par\vspace{0.4em}}

% --- Writelines helper: blank ruled lines for short-answer responses ---
\newcounter{writelinesi}
\newcommand{\writelines}[1]{%
  \par\vspace{4pt}%
  \setcounter{writelinesi}{0}%
  \loop\ifnum\value{writelinesi}<#1
    \par\noindent\rule{\linewidth}{0.4pt}\vspace{1.4em}%
    \stepcounter{writelinesi}%
  \repeat
  \par\vspace{4pt}}

% --- List metrics: editorial proportions, tightened ---
\setlist[enumerate,1]{leftmargin=2.4em, itemsep=0.45em, topsep=0.45em, parsep=0pt, label=\textbf{\arabic*.}}
\setlist[enumerate,2]{leftmargin=1.8em, itemsep=0pt, topsep=0.1em, parsep=0pt, label=(\alph*)}

% --- Code block style ---
\lstset{
  basicstyle=\ttfamily\small,
  breaklines=true,
  breakatwhitespace=true,
  columns=fullflexible,
  showstringspaces=false,
  numbers=none,
  tabsize=3,
  aboveskip=4mm,
  belowskip=4mm,
  frame=tb,
  framerule=0.4pt,
  framesep=6pt,
  xleftmargin=12pt,
  xrightmargin=12pt,
}

% --- Even-page padding for duplex printing ---
\AtEndDocument{%
  \clearpage%
  \ifodd\value{page}\else%
    \thispagestyle{plain}%
    \null\vfill%
    \begin{center}\textit{\small This page intentionally left blank.}\end{center}%
    \vfill\null%
  \fi%
}

\setlength{\emergencystretch}{2em}
\sloppy

% --- Key banner activation (answers mode) ---
\ifanswers\fancyhead[R]{\textit{\small Exam Name --- \textcolor{keyred}{\textbf{KEY}}}}\fi

\date{}

\begin{document}

\metastrip{CECS NNN}{Term YYYY}{Exam Name}{100}   % <-- edit all four fields
\par\vspace{0.4em}
\begin{center}{\Large\itshape Exam Name}\end{center}   % <-- edit exam name
\par\vspace{-0.2em}
\doublerule
\par\vspace{0.9em}

\par\vspace{6pt}
\noindent\identityinstruction
\par\vspace{12pt}

\noindent{\sffamily\bfseries\large NAME:}~\fieldline{13.5cm}{\large\studentname}
\par\vspace{14pt}
\noindent{\sffamily\bfseries\large STUDENT ID\#:}~\fieldline{5.5cm}{\large\ttfamily\studentid}\hfill{\sffamily\bfseries\large DATE:}~\rule[-3pt]{4cm}{0.8pt}
\par\vspace{6pt}

\ifanswers\par\vspace{0.4em}\noindent\colorbox{keyred}{\parbox{\dimexpr\textwidth-2\fboxsep\relax}{\centering\color{white}\sffamily\bfseries\large $\star$~~ANSWER KEY~~---~~NOT FOR DISTRIBUTION~~$\star$}}\par\vspace{0.4em}\fi

\par\vspace{0.6em}
\noindent\ding{43}~\textbf{Directions.}~\textit{Write on the exam paper. You may use one 3''$\times$5'' card of notes (both sides, handwritten only). \textbf{No electronic devices of any kind are allowed on your desk during the exam.}}
\par\vspace{0.4em}

% ============================================================
%  QUESTION BODY — replace the placeholders below with your
%  actual questions. Keep \examsection, \correctchoice /
%  \wrongchoice, and \ifanswers gating intact.
% ============================================================

% --- Section 1: Multiple Choice (2 pts each) ---
\examsection{Multiple Choice}
\begin{enumerate}

\item \textit{(2 pts)}~Placeholder MC stem (replace per exam).
  \begin{enumerate}[label=(\alph*)]
    \item \correctchoice{option A}
    \item \wrongchoice{option B}
    \item \wrongchoice{option C}
    \item \wrongchoice{option D}
  \end{enumerate}
  \ifanswers \textcolor{keyred}{\textbf{Answer:} a} \fi

\end{enumerate}

% --- Section 1b: True / False ---
%  House standard: the \textsc{T~/~F.} label keeps the question typed `tf`
%  (parse_outline_from_tex / _classify), and the (a) True / (b) False choices
%  MUST be a stacked enumerate — one option per line — so Gradescope's region
%  detection can find them. Do NOT inline the options. Always (a) True, (b)
%  False; wrap the correct one in \correctchoice. The inline Answer: line is the
%  canonical record parse_outline_from_tex reads (keep it "True"/"False").
\examsection{True / False}
\begin{enumerate}[resume]

\item \textit{(2 pts)}~\textsc{T~/~F.}~Placeholder true/false claim (replace per exam).
  \begin{enumerate}[label=(\alph*)]
    \item \correctchoice{True}
    \item \wrongchoice{False}
  \end{enumerate}
  \ifanswers \textcolor{keyred}{\textbf{Answer:} True} \fi

\end{enumerate}

% --- Section 2: Short Answer (model answer behind \ifanswers) ---
\examsection{Short Answer}
\begin{enumerate}[resume]

\item \textit{(10 pts)}~Placeholder short-answer prompt (replace per exam).
  \ifanswers
    \par\smallskip\textit{Key:} model answer here.
  \else
    \writelines{6}
  \fi

\end{enumerate}

\par\vfill
\begin{center}\small\textit{End of exam.}\end{center}

\end{document}
